% !TeX document-id = {f5eab2f7-1326-43b3-ac0e-350c2acbded8}

% !TeX TXS-program:compile = txs:///pdflatex/[--shell-escape]

\documentclass[a4paper, fontsize=11pt, parskip=half, twoside]{scrreprt}

\usepackage[utf8]{inputenc}
\usepackage[T1]{fontenc}   
\usepackage{graphicx}       
\usepackage[english, ngerman]{babel}
% for displaying quotes
\usepackage{csquotes}     
\usepackage{acronym}
\usepackage{eurosym}
\usepackage[linktocpage=true]{hyperref}
\usepackage{xurl}
\usepackage[bindingoffset=8mm]{geometry}
\usepackage{caption}
\captionsetup{format=hang, justification=raggedright}
\usepackage[style=authoryear, backend=biber]{biblatex}
\usepackage{float}
\usepackage{rotating}
\usepackage{amsmath}
\usepackage{amssymb}

% for nested lists
\usepackage{outlines}

% for drawing
\usepackage{tikz}

% improve micro typography for better distribution
\usepackage{microtype}

% adds generic commands for degree, ohm, ..
\usepackage{textcomp}
\usepackage{gensymb}

% for source code
\usepackage{minted}

% caption for e.g. figures
\usepackage{caption}
% captions for multiple figures
\usepackage{subcaption}
% for custom dates
\usepackage{datetime}
% clever refs
% set labels as always (\label{fig:hello}), \cref{fig:hello}
\usepackage[german]{cleveref}

% 1.5x line space
\usepackage[onehalfspacing]{setspace}

\usepackage{todonotes}

% define date variable for whole doc
% use \displaydate{date} where needed to insert this date

\newdate{date}{01}{07}{2023}
\date{\displaydate{date}}

% add zotero file for citations
\addbibresource{reference.bib}

\begin{document}
	
	\thispagestyle{empty}
	
	\cleardoublepage   % force output to a right page
	\thispagestyle{empty}
	\begin{titlepage}
		\begin{flushright}
			\includegraphics[width=0.4\linewidth]{assets/Logo-A3.jpg}
		\end{flushright}
		
		\begin{flushleft}
			\section*{Integration eines Feature-orientierten Testsystems in den Entwicklungszyklus technischer Systeme}
			%\subsection*{\papersubtitle}
			\vspace{1cm}
			
			\vspace{0.5cm}
			Bachelorarbeit\newline
			zum Erlangen des akademischen Grades\newline
			
			\vspace{0.5cm}
			\textbf{Bachelor of Science in Engineering (BSc)}
			
			\vspace{1cm}
			Fachhochschule Vorarlberg\newline
			Informatik - Software and Information Engineering\newline
			
			\vspace{0.5cm}
			Betreut von\newline
			Dipl.-Ing. Dr. techn. Ralph Hoch
			
			\vspace{0.5cm} 
			Vorgelegt von\newline
			Marco Prescher	
			
			\vspace{0.5cm}
			Dornbirn, am \displaydate{date}
		\end{flushleft}
	\end{titlepage}
	
	% Kurzreferat:
	\newpage
	\section*{Kurzreferat}
	\subsection*{Integration eines Feature-orientierten Testsystems in den Entwicklungszyklus technischer Systeme}
	
	Die Entwicklung technischer Systeme ist ein komplexer und kostspieliger Prozess. 
	Daher ist es wichtig, dass die Produkte vor der Auslieferung an den Kunden gezielt und sorgfältig getestet werden. 
	Dies wird durch Zeitdruck und Deadlines oftmals vernachlässigt, oder nur unzureichend durchgeführt. 
	Mit einem gut strukturierten Testplan kann dieses Risiko allerdings minimiert werden, und die Durchführung der Tests mit dem Produktentwicklungszyklus integriert werden. 
	Dadurch kann stabile Hardware sowie effiziente und gut strukturierte Software ohne große Verzögerung an den Kunden ausgeliefert werden.
	
	Durch Methodiken wie zum Beispiel Feature-orientierte Entwicklung ist es möglich, dass bestimmte Features vor dem Release der Produkte abgenommen und getestet werden müssen. 
	Das wiederum ermöglicht, dass neu implementierte Features gründlich getestet werden und somit zu einer hohen Qualität beitragen.
	
	Dies kann erreicht werden, indem Features strukturiert in einer Datenbank eingepflegt werden. 
	Ein auf diesen Featurebeschreibungen basierendes System kann die automatische Testdurchführung unterstützen, gezielt Tests für einzelne Features durchführen und deren Abnahme beschleunigen. 
	Dadurch kann der Testaufwand verringert und den Entwicklern ein fokussiertes Feedback vermittelt werden.
	
	\newpage
	\section*{Abstract}
	\subsection*{Integrating a feature-oriented testing system into the development cycle of technical systems}
	
	The development of technical systems is a complex and costly process. 
	Therefore, it is important that products are thoroughly tested before delivery to the customer. 
	However, this is often neglected or done insufficiently due to time pressure and deadlines. 
	A well-structured test plan can minimize this risk and integrate the testing process into the product development cycle. 
	This allows stable hardware and efficient, well-structured software to be delivered to the customer without significant delays.
	
	Using techniques such as feature-oriented development, certain features must be accepted and tested before the product release. 
	This in turn allows newly implemented features to be thoroughly tested and contribute to high quality. 
	
	This can be achieved by structuring features in a database. 
	A system based on these feature descriptions can support automatic testing, conduct targeted tests for individual features, and accelerate their acceptance. 
	This reduces testing effort and provides focused feedback to the developers.
	
	% Inhaltsverzeichnis:
	\cleardoublepage   % force output to a right page
	\setcounter{tocdepth}{2}
	\setcounter{secnumdepth}{4}
	\tableofcontents
	
	% Abbildungsverzeichnis
	\clearpage
	\phantomsection
	\addcontentsline{toc}{chapter}{Abbildungsverzeichnis}
	\listoffigures
	
	% Abkürzungsverzeichnis:
	\clearpage
	\phantomsection
	\addcontentsline{toc}{chapter}{Abkürzungsverzeichnis}
	\section*{Abkürzungsverzeichnis}
	\begin{acronym}
		\acro{API}{Application Programming Interface}
		\acro{JSON}{JavaScript Object Notation}
		\acro{TCMS}{Test Case Management System}
		
		\acrodefplural{TCMS}{Test Case Management Systeme}
	\end{acronym}
	
	\clearpage
	\section*{Danksagung}
	Ich möchte mich aufrichtig bei Ralph Hoch, der Firma Gantner Instruments und allen Mitarbeitenden bedanken, die mir bei der Vollendung dieser Arbeit geholfen haben. 
	Ihre Unterstützung und Expertise waren von unschätzbarem Wert und haben zum erfolgreichen Abschluss dieses Projekts beigetragen. 
	Vielen Dank für Ihre harte Arbeit und Ihr Engagement. 
	Ich schätze Ihre Zusammenarbeit sehr.
	
	
	
	\clearpage
	\chapter{Einleitung}
	\todo{Proper Einleitung Description}
	
	
	\section{Motivation} \label{sec:motivation}
	Die Entwicklung technischer Systeme ist ein komplexer Prozess, der eine hohe Qualität erfordert, um den Anforderungen der Kunden gerecht zu werden. 
	Damit diese Qualität auch gewährleistet wird, müssen Fehler sowie Mängel identifiziert und ausgebessert werden. 
	Ein \ac{TCMS} bietet eine Lösung, um diesen Prozess zu vereinfachen und effektiver zu gestalten. 
	Durch die Verwendung eines \ac{TCMS} können Angestellte aus verschiedenen Abteilungen, wie z.B. der Software-, Hardware-, Support- oder Marketingabteilung, die Qualität des Produkts gemeinsam verbessern. 
	Diese Arbeit fokussiert sich daher auf die Integration von einem \ac{TCMS} in einen laufenden Produktentwicklungs-Zyklus und vergleicht verschiedene vorhandene Lösungen.
	
	\section{Problemstellung}
	Durch die von \autoref{sec:motivation} angesprochene Komplexität technischer Systeme ist bekannt, dass Fehler, die erst spät im Entwicklungsprozess entdeckt werden, viel kostspieliger zu beheben sind als Fehler, die frühzeitig identifiziert und behoben werden. 
	Infolgedessen suchen Unternehmen nach Lösungen, um den Testprozess effektiver zu gestalten und Fehler früher im Entwicklungszyklus zu identifizieren. 
	\ac{TCMS} bietet eine solche Lösung, indem es Entwicklern und Testern ermöglicht, Testfälle effizient zu planen und zu verwalten sowie Testergebnisse zu erfassen, darzustellen und somit auch zu verfolgen. 
	Obwohl \aclp{TCMS} in der Industrie weit verbreitet sind und es einige fertige Lösungen gibt, gibt es jedoch nicht immer die Perfekte Lösung um ein bestehendes \ac{TCMS} für das eigene Projekt anzuwenden. 
	Insbesondere gibt es Bedenken hinsichtlich der Anwendbarkeit von \ac{TCMS} mit anderen Tools und der Skalierbarkeit von \ac{TCMS}. 
	Diese Probleme stellen Hindernisse dar, die die Einführung von \ac{TCMS} in einem Unternehmen erschweren können. 
	In dieser Arbeit werden wir uns mit diesen Problemen und Herausforderungen auseinandersetzen, eine Software Lösung entwickeln und untersuchen, wie \aclp{TCMS} effektiv eingesetzt werden können, um den Testprozess zu optimieren und die Qualität eines Produkts zu verbessern.
	
	\section{Zielsetzung}
	Das Ziel dieser Arbeit ist es, die Verwendung von \aclp{TCMS}n zu untersuchen und zu bewerten. 
	Zudem ein auf unser eigenes Produkt angepasstes \ac{TCMS} backend zu entwickeln. 
	Hiermit ergeben sich drei relevante fragen:
	
	\begin{itemize}
		\item Wie kann man Features, Testfälle und Testimplementierungen beschreiben? 
		\item Wie in Datenbank schreiben und lesen? 
		\item Wie mit Testläufen verknüpfen?
	\end{itemize}
	
	Insbesondere möchten wir die folgenden Ziele erreichen:
	
	\begin{itemize}
		\item Die Vor- und Nachteile der Verwendung von \ac{TCMS} zu identifizieren und zu analysieren.
		\item Die Entwicklung und Integration von \ac{TCMS} in den Entwicklungsprozess zu untersuchen und zu bewerten, einschließlich der Herausforderungen, die bei der Integration von \ac{TCMS} in bestehende Entwicklungs- und Testprozesse auftreten können.
		\item Die Konnektivität von \ac{TCMS} mit anderen Tools und Systemen, die in der Softwareentwicklung verwendet werden, zu untersuchen und zu bewerten.
		\item Empfehlungen für die erfolgreiche Implementierung von \ac{TCMS} in der Softwareentwicklung zu geben, einschließlich der Identifizierung bewährter Praktiken.
	\end{itemize}
	
	Durch die Erfüllung dieser Ziele wird diese Arbeit dazu beitragen, das Verständnis für die Verwendung von \ac{TCMS} in der Produktentwicklung zu verbessern und Unternehmen dabei zu unterstützen, den Testprozess zu optimieren und die Qualität ihrer Produkte zu verbessern.
	
	
	
	\chapter{Stand des Wissens}
	Dieses Kapitel gibt einen Überblick über den aktuellen Stand der Technik von Testarten und Testmanagementsystemen. 
	Darüber hinaus wird untersucht wie ein \acl{TCMS} verwendet wird und was die empfohlene Vorgehensweise ist.
	Zusätzlich werden Technologien analysiert, mit denen eine Implementation von einem \ac{TCMS} gut umsetzbar ist. 
	
	
	\section{Testarten} \label{sec:testtypes}
	Generell gibt es einige Testarten in der Entwicklung eines Produktes. 
	Es gibt laut \textcite{atlassian_unterschiedlichen_nodate} mindestens sieben unterschiedliche Testverfahren die je nachdem andere Implementierungen betrachten und testen.
	In \textcite{noauthor_software_nodate} wird beschreiben, dass dabei zwischen Funktionalen und Nicht-Funktionalen Tests unterschieden werden.
	Zu der funktionalen Testfamilie zählen beispielsweise Unit-Tests und Integrationstests.
	Wobei hingegen Leistung, Last und Stresstests sowie Usability-Tests zu den nicht funktionalen Testfamilie gehören.
	
	Ein paar der meist vorkommenden Testarten sind:
	
	\begin{itemize}
		\item Unit-Tests
		\item Integrationstests
		\item Funktionstests
		\item Leistungstests
	\end{itemize}

	Um eine Testabdeckung sicherzustellen und den allgemeinen Testprozess zu optimieren, gibt es, wie in \autoref{sec:motivation} auf \autopageref{sec:motivation} schon besprochen, \aclp{TCMS}. 
	Diese Systeme bieten eine zentrale Plattform zur Verwaltung von Testfällen, dass alle erforderlichen Tests durchgeführt und dokumentiert werden. 
	In \autoref{sec:tcms} auf \autopageref{sec:tcms} werden wir diese Systeme genauer untersuchen.
	
	\subsection{Unit-Tests}
	Im Allgemeinen sind Unit-Tests sehr einfache Tests die beispielsweise Methoden einer Klasse mit unterschiedlichen Parametern testet. 
	Sie sind automatisierbar und können von einer Continuos-Integration-Pipeline sehr schnell durchgeführt werden. 
	Dieser Ansatz ermöglicht eine frühe Entdeckung von Fehlern und eine kontinuierliche Überprüfung der Funktionsfähigkeit des Produkts im Entwicklungsprozess.
	
	\subsection{Integrationstests}
	Integrationstests sind Tests, die sicherzustellen, dass verschiedene Module oder Services, problemlos miteinander interagieren können. 
	Durch diese Tests kann die Funktionalität der Anwendung insgesamt überprüft werden, wie beispielsweise die Interaktion mit einer Datenbank oder der Zusammenarbeit von Mikroservices.
	
	\subsection{Funktionstests}
	Funktionstests werden integriert, um ausschließlich Ergebnisse einer gegebenen Funktion zu überprüfen. Dabei werden Spezifikation vor der Testausführung festgelegt und diese dann mit den Ergebnissen verglichen.
	
	Im Gegensatz dazu stellt ein Integrationstest sicher, dass verschiedene Komponenten einer Anwendung zusammenarbeiten.
	Dabei werden Interaktionen zwischen verschiedenen Modulen oder Services, wie z.B. Datenbankabfragen oder die Kommunikation zwischen Mikroservices, überprüft.
	
	Während Integrationstests beispielsweise nur prüfen, ob Datenbankabfragen generell möglich sind, wird bei einem Funktionstest ein bestimmter Wert aus der Datenbank abgerufen und dieser mit den angegebenen Spezifikationen geprüft.
	
	\subsection{Leistungstests}
	Leistungstests sind Verfahren, die dazu dienen, das Verhalten eines Systems unter verschiedenen Auslastungsmaximierungen zu überprüfen. 
	Sie werden oft dazu verwendet um Zuverlässigkeit, Geschwindigkeit, Skalierbarkeit und Reaktionsfähigkeit einer Anwendung zu testen. 
	Außerdem können Leistungstests mögliche Engpässe in einer Anwendung identifizieren.
	
	\section{Blackbox/Whitebox Testing}
	\todo{Proper Blackbox/Whitebox Testing Description}
	
	
	\section{Test Driven Devlopment}
	\todo{Proper Test Driven Devlopment Description}
	
	
	\section{Test Case Management System} \label{sec:tcms}
	Ein \acl{TCMS} ist eine Software, mit dem Test-Teams für ein bestimmtes Projekt oder eine Anwendung Testfälle verwalten, organisieren und verfolgen können.
	Es hilft bei der Planung, Überwachung und Dokumentation von Tests und ermöglicht es, Testfälle sicher und effizient zu verwalten.
	
	Zusammenfassend ist das \ac{TCMS} ein wichtiges Werkzeug, um einen strukturierten und effektiven Testprozess zu gewährleisten und den Qualitätsstandard einer Anwendung zu verbessern.

	\section{Merkmale von Test Case Management Systemen} \label{sec:functionsOfTCMS}
	Dieser Abschnitt gibt einen Überblick zu den Hauptmerkmalen von \aclp{TCMS}.
	Ein \ac{TCMS} verfügt normalerweise über Basisfunktionen wie Testfall-Erstellung, Testfall-Verwaltung, Testfall-Ausführung und Ergebnisberichterstattung. 
	Es kann auch eine integrierte Umgebung für die Zusammenarbeit von Testern und Entwicklern bereitstellen.
	
	Zusätzlich zu den Basisfunktionen sind das die 6 wichtigsten Merkmale, die eine gutes \ac{TCMS} ausmachen:
	
	\begin{itemize}
		\item Attraktive Benutzeroberfläche und benutzerfreundliches Design
		\item Zulassung von mehreren Projekten und Benutzerberechtigungen
		\item Nachvollziehbarkeit
		\item einfache Zeitplanung und Organisation
		\item Überwachung und Metriken
		\item Flexibilität
	\end{itemize}
	
	\textcite{lead_articles_nodate}	
	
	\subsection{Attraktive Benutzeroberfläche und benutzerfreundliches Design}
	\todo{Proper Attraktive Benutzeroberfläche und benutzerfreundliches Design Description}
	\textcite{lead_articles_nodate}	
	
	\subsection{Zulassung von mehreren Projekten und Benutzerberechtigungen}
	\todo{Proper Zulassung von mehreren Projekten und Benutzerberechtigungen Description}
	\textcite{lead_articles_nodate}	
	
	\subsection{Nachvollziehbarkeit}
	\todo{Proper Nachvollziehbarkeit Description}
	\textcite{lead_articles_nodate}	
	
	\subsection{Zeitplanung und Organisation}
	\todo{Proper Zeitplanung und Organisation Description}
	\textcite{lead_articles_nodate}
	
	\subsection{Überwachung und Metriken}
	\todo{Proper Überwachung und Metriken Description}
	\textcite{lead_articles_nodate}	
	
	\subsection{Flexibilität}
	\todo{Proper Flexibilität Description}
	\textcite{lead_articles_nodate}	
	
	
	\section{Allgemeine Vorteile von Test Case Management Systemen}
	Der Hauptvorteil des Einsatzes von einem \ac{TCMS} besteht darin, dass sie den Testprozess verbessern und so schneller zu einem besseren Produkt führen. 
	Sie sind auch unerlässlich, um die Gesamtkosten des Testens zu kontrollieren, indem sie die Testautomatisierung nutzen, um einen reibungslosen Ablauf zu gewährleisten.  
	
	Im Folgenden werden die spezifischen Vorteile von Testmanagement-Tools für die Testausführungsabläufe erläutert:
	
	\begin{itemize}
		\item Sie geben einen besseren Überblick über das zu testende System, halten den gesamten Prozess auf Kurs und koordinieren die Testaktivitäten.
		\item Sie helfen bei der Feinabstimmung des Testprozesses, indem sie die Zusammenarbeit, Kommunikation und Datenanalyse unterstützen.
		\item Sie verfolgen Aufgaben, Fehler und Testergebnisse und vereinfachen den Prozess, indem sie alles in einer einzigen Anwendung erledigen.
		\item Sie sind skalierbar und können eingesetzt werden, wenn die Testaktivitäten umfangreicher und komplexer werden.
	\end{itemize}

	\textcite{lead_articles_nodate}
	
	\section{Überblick von bestehenden Test Case Management Systeme}
	Heutzutage gibt es natürlich schon eine Vielzahl von fertigen und direkt verwendbaren \aclp{TCMS}n auf dem Markt, die die Verwaltung von Tests vereinfacht.
	
	Einige bekannte Tools sind:
	
	\begin{itemize}
		\item \textcite{noauthor_testrail_2023} 
		
		Webbasiertes \acl{TCMS}, ist zentralisiert und hat eine einfache und intuitive Oberfläche.
		
		Features:
		\begin{itemize}
			\setlength\itemsep{-0.5em}
			\item Test Planung, Verwaltung und Ausführung 
			\item Echtzeit Berichterstattung und Analysen
			\item Rückverfolgbarkeits- und Abdeckungsberichte für Anforderungen, Tests und Fehler
		\end{itemize}
		
		\item \textcite{noauthor_tricentis_nodate} 
		
		Webbasiertes \acl{TCMS} und hat moderne und intuitive Oberfläche.
		
		Features:
		\begin{itemize}
			\setlength\itemsep{-0.5em}
			\item Test Planung, Verwaltung und Ausführung 
			\item Echtzeit Test Status Update
			\item Berichterstattung
		\end{itemize}
		
		\item \textcite{noauthor_zephyr_nodate} 
		
		Webbasiertes \acl{TCMS}, moderne Oberfläche und kann in JIRA integriert werden. 
		
		Features:
		\begin{itemize}
			\setlength\itemsep{-0.5em}
			\item Test Planung, Verwaltung und Ausführung 
			\item Test Automatisierung
			\item Echtzeit Visualisierung von Projektstatus
		\end{itemize}
		
		\item \textcite{noauthor_tricentis_nodate-1} 
		
		Webbasiertes \acl{TCMS} und hat moderne und intuitive Oberfläche.
		
		Features:
		\begin{itemize}
			\setlength\itemsep{-0.5em}
			\item Test Planung, Verwaltung und Ausführung 
			\item Migrationsmöglichkeit von alten Test Management Lösungen
			\item Integrationsmöglichkeit mit Jenkins, Azure Pipelines, Bamboo oder jedem anderen CI/CD-Tool
			\item Anpassbare Dashboards, um über alle Releases, Projekte oder Programme im gesamten Unternehmen zu berichten
			\item Berichte per E-Mail oder URL teilen
		\end{itemize}
	\end{itemize}
	
	Diese \aclp{TCMS} bieten eine Grundfunktionalität, auf diese wir in \autoref{sec:functionsOfTCMS} auf \autopageref{sec:functionsOfTCMS} eingegangen sind. 
	Zu den Grundfunktionalitäten bieten diese Systeme jeweils andere oder zusätzliche Funktionen und Integrationsmöglichkeiten.
	Welches Tool am besten geeignet ist, hängt natürlich von der jeweiligen Anwendung ab.
	
	\begin{figure}[H]
		\centering
		\includegraphics[scale=0.5]{assets/best-tool-logos.jpg}
		\caption{Liste von aktuell führenden Test Case Management Systemen/Tools (Quelle: \textcite{noauthor_7_nodate})}
		\label{fig:tcms_logo}
	\end{figure}
	
	Die in Abbildung \ref{fig:tcms_logo} ersichtlichen Tools zeigen eine aktuell führende Auswahl von beliebten \aclp{TCMS}n.  
	
	\chapter{Konzept}
	\todo{Proper Lösungsansatz}
	
	
	\section{Scrum}
	\todo{Proper Scrum Description}
	
	
	\section{Architektur}
	\todo{Proper Architektur Description}
	
	\subsection{Backend}
	\todo{Proper Backend Description}
	
	\subsection{Domain Driven Design}
	\todo{Proper Domain Driven Design Description}
	
	
	\section{Vorentwicklungsphase}
	\todo{Proper Vorentwicklungsphase Description}
	
	\subsection{Projektstruktur}
	\todo{Proper Projektstruktur Description}
	
	\subsection{Domain Model}
	\todo{Proper Domain Model Description}
	
	\subsection{ER Model}
	\todo{Proper ER Model Description}
	
	\subsection{Featurestruktur}
	\todo{Proper Featurestruktur Description}
	
	\subsection{Featurebedingungen}
	\todo{Proper Featurebedingungen Description}
	
	\subsection{Featureverwaltung}
	\todo{Proper Featureverwaltung Description}
	
	
	\section{Entwicklungsphase}
	\todo{Proper Entwicklungsphase Description}
	
	\subsection{Verwaltung und Struktur von Features}
	\todo{Proper Verwaltung und Struktur von Features Description}
	
	\subsection{Featurestruktur zur Datenbank}
	\todo{Proper Featurestruktur zur Datenbank Description}
	
	\subsection{Visualisierung von Featurestrukturen}
	\todo{Proper Visualisierung von Featurestrukturen Description}
	
	
	
	\chapter{Technische Umsetzung}
	\todo{Proper Implementierung}
	
	
	\section{Verwendete Technologien}
	Dieser Abschnitt gibt eine Übersicht über die verwendeten Technologien, die für die implementieren, von einem backend eines \acl{TCMS}s eingesetzt worden sind.
	
	\subsection{MariaDB}
	\todo{Proper MariaDB Description}
	\textcite{noauthor_mariadb_nodate}
	
	\subsection{Docker}
	\todo{Proper Docker Description}
	
	\subsection{Microsoft .NET Core}
	\todo{Proper .NET Core Description}
	\textcite{billwagner_net_nodate}
	
	\subsection{Microsoft ASP.NET Core}
	\todo{Proper ASP.NET Core Description}
	\textcite{billwagner_net_nodate}
	
	\subsection{Entity Framework Core}
	\todo{Proper Entity Framework Core Description}
	\textcite{billwagner_net_nodate}
	
	\subsection{OpenAPI}
	\todo{Proper OpenAPI Description}
	\textcite{noauthor_openapi_nodate}
	\textcite{noauthor_openapi_nodate-1}
	
	\section{Sprints}
	\todo{Proper Sprints Description}
	\textcite{noauthor_scrum_nodate}
	
	\section{User Stories}
	\todo{Proper User Stories Description}
	
	\subsection{Story 1}
	\todo{Proper Story 1 Description}
	
	\subsection{Story 2}
	\todo{Proper Story 2 Description}
	
	\subsection{Story n}
	\todo{Proper Story n Description}
	
	
	
	\chapter{Ergebnisse}
	\todo{Proper Ergebnisse Description}
	
	
	
	\chapter{Fazit und Ausblick}
	\todo{Proper Fazit und Ausblick}
	
	
	
	\begin{figure}[ht]
		\begin{minted}[autogobble, tabsize=4]{rust}
			{
				// Start eines neuen Scopes
				
				let hello = String::from("hello");
				// Ende des Scopes. hello wird an dieser Stelle freigegeben
			}
			// hello kann hier nicht mehr verwendet werden
		\end{minted}
		\caption{Ein String wird in einem Scope erzeugt und der Variable hello zugewiesen. Am Ende des Scopes wird der Speicher freigegeben.}
		\label{fig:drop:example}
	\end{figure}
	
	\begin{figure}[H]
		\centering
		\includegraphics[width = 10cm]{assets/Logo-A3.jpg}
		\caption{Nach 100 Simulationsschritten, einem minimalen Gruppenanteil $B{\min} = 0.4$ und einer Nachbarschaftsgröße von einem Feld sind die Gruppen bereits sichtlich separiert.}
		\label{fig:schelling_end}
	\end{figure}
	
	% Literaturverzeichnis:
	\clearpage
	\phantomsection
	\addcontentsline{toc}{chapter}{Literaturverzeichnis}
	\printbibliography
	
	\clearpage
	\section*{Eidesstattliche Erklärung}
	Ich erkläre hiermit an Eides statt, dass ich vorliegende Bachelorarbeit selbstständig und ohne Benutzung anderer als der angegebenen Hilfsmittel angefertigt habe. 
	Die aus fremden Quellen direkt oder indirekt übernommenen Stellen sind als solche kenntlich gemacht. 
	Die Arbeit wurde bisher weder in gleicher noch in ähnlicher Form einer anderen Prüfungsbehörde vorgelegt und auch noch nicht veröffentlicht.
	
	\vspace{3cm}
	\noindent
	Dornbirn, am \displaydate{date} \hfill Marco Prescher
	
	
\end{document}